% =========================================================
%  sections/08-penilaian.tex
%  Bagian 8 — Kriteria, Indikator, dan Bobot Penilaian
% =========================================================

\section{Kriteria, Indikator, dan Bobot Penilaian}
\label{sec:penilaian}

\subsection{Prinsip Penilaian}

\begin{itemize}
  \item Penilaian berbasis kriteria (\istilah{criterion-referenced}), bukan
        berdasarkan sebaran normal nilai kelas. Mahasiswa dinilai terhadap
        kriteria ketercapaian Sub-CPMK yang telah ditetapkan, bukan terhadap
        prestasi relatif teman sekelas.
  \item Seluruh teknik dan alat asesmen dipetakan secara eksplisit ke Sub-CPMK
        dan CPMK, sehingga setiap nilai yang diperoleh dapat ditelusuri
        kontribusinya terhadap CPL.
  \item Bobot penilaian ditentukan sesuai kontribusi komponen terhadap
        pencapaian CPMK, bukan berdasarkan kemudahan pelaksanaannya.
  \item Penilaian bersifat autentik: mahasiswa dinilai melalui artefak nyata
        berupa kode program, level, \istilah{build}, dan laporan, bukan semata
        melalui hafalan.
  \item Umpan balik diberikan secara formatif pada setiap praktikum agar
        mahasiswa dapat memperbaiki pekerjaannya sebelum asesmen sumatif.
\end{itemize}

\subsection{Komponen Penilaian dan Bobot}

\begin{table}[htbp]
\centering
\caption{Komponen penilaian, instrumen, dan bobot}
\label{tab:komponen-nilai}
\begin{tabularx}{\textwidth}{@{}L{3.6cm} Y L{2.5cm} C{1.6cm}@{}}
\toprule
\headrow \textbf{Komponen} & \textbf{Bentuk Instrumen} & \textbf{CPMK yang Diukur} & \textbf{Bobot} \\
\midrule
Tugas Praktikum dan Portofolio &
Penilaian artefak kode mingguan melalui repositori Git; enam sampel dinilai menggunakan rubrik kesesuaian fungsional, kualitas kode, dan kelengkapan dokumentasi. &
\CPMK{1}--\CPMK{4} & 25\% \\
\addlinespace
Kuis &
Tes singkat dua kali (15--20 menit) untuk mengukur pemahaman konsep secara formatif. &
\CPMK{1}, \CPMK{2} & 10\% \\
\addlinespace
Ujian Tengah Semester (UTS) &
Tes tulis konsep dan uji kode (\istilah{debugging} serta penyempurnaan kode Pygame dalam batas waktu). &
\CPMK{1}, \CPMK{2} & 20\% \\
\addlinespace
Mini Proyek AI &
Penilaian produk kelompok (kode FSM dan A*), laporan \istilah{playtest}, dan presentasi. &
\CPMK{2} & 10\% \\
\addlinespace
Proyek Akhir &
Penilaian prototipe Unity (level, aset visual, audio, UI), laporan pengembangan, presentasi, dan uji fungsional pada UAS. &
\CPMK{3}, \CPMK{4} & 30\% \\
\addlinespace
Partisipasi dan Kehadiran &
Observasi kehadiran dan keaktifan dalam praktikum, diskusi, serta kegiatan \istilah{playtest}. &
\CPMK{1}--\CPMK{4} & 5\% \\
\midrule
\textbf{Total} & & & \textbf{100\%} \\
\bottomrule
\end{tabularx}
\end{table}

\subsection{Matriks Bobot Komponen terhadap CPMK}

\begin{table}[htbp]
\centering
\caption{Distribusi bobot setiap komponen penilaian pada CPMK (dalam persen nilai akhir)}
\label{tab:matriks-bobot}
\begin{tabularx}{\textwidth}{@{}L{4.6cm} C{1.8cm} C{1.7cm} C{1.7cm} C{1.7cm} C{1.7cm}@{}}
\toprule
\headrow \textbf{Komponen} & \textbf{Bobot} & \textbf{\CPMK{1}} & \textbf{\CPMK{2}} & \textbf{\CPMK{3}} & \textbf{\CPMK{4}} \\
\midrule
Tugas Praktikum dan Portofolio & 25\% & 8 & 6 & 6 & 5 \\
Kuis                           & 10\% & 5 & 5 & --- & --- \\
Ujian Tengah Semester          & 20\% & 10 & 10 & --- & --- \\
Mini Proyek AI                 & 10\% & --- & 10 & --- & --- \\
Proyek Akhir                   & 30\% & --- & --- & 15 & 15 \\
Partisipasi dan Kehadiran      & 5\%  & 2 & 1 & 1 & 1 \\
\midrule
\textbf{Total per CPMK}        & \textbf{100\%} & \textbf{25\%} & \textbf{32\%} & \textbf{22\%} & \textbf{21\%} \\
\bottomrule
\end{tabularx}
\end{table}

\subsection{Kontribusi CPMK terhadap CPL}

\begin{table}[htbp]
\centering
\caption{Kontribusi CPMK terhadap pencapaian CPL pada mata kuliah ini}
\label{tab:kontribusi-cpl}
\begin{tabularx}{\textwidth}{@{}L{2.4cm} L{6.4cm} C{2.4cm} Y@{}}
\toprule
\headrow \textbf{CPL} & \textbf{CPMK Penyusun} & \textbf{Bobot} & \textbf{Bukti Ketercapaian} \\
\midrule
\CPL{03} & \CPMK{1} (25\%) dan \CPMK{2} (32\%) & 57\% &
Kode program Pygame mahasiswa, hasil kuis, dan jawaban UTS yang menunjukkan kemampuan merancang serta menganalisis algoritma permainan. \\
\addlinespace
\CPL{04} & \CPMK{3} (22\%) dan \CPMK{4} (21\%) & 43\% &
Level Unity, prototipe akhir beserta asetnya, laporan pengembangan, dan hasil uji fungsional serta \istilah{playtest}. \\
\midrule
\textbf{Total} & & \textbf{100\%} & \\
\bottomrule
\end{tabularx}
\end{table}

\subsection{Rubrik Penilaian}

\subsubsection{Rubrik Mini Proyek AI}

\begin{table}[htbp]
\centering
\caption{Rubrik penilaian mini proyek AI}
\label{tab:rubrik-miniproyek}
\begin{tabularx}{\textwidth}{@{}L{4.4cm} C{1.6cm} Y@{}}
\toprule
\headrow \textbf{Kriteria} & \textbf{Bobot} & \textbf{Deskripsi Capaian} \\
\midrule
Implementasi FSM & 25\% &
NPC memiliki \istilah{state} yang jelas dan transisi yang berjalan sesuai kondisi jarak serta \istilah{cooldown}; tidak terjadi kebuntuan \istilah{state}. \\
\addlinespace
Implementasi A* & 25\% &
A* menghasilkan jalur yang benar pada peta berhalangan, menggunakan struktur data antrean prioritas, dan heuristik yang dapat dipertanggungjawabkan. \\
\addlinespace
Integrasi FSM dan A* & 20\% &
Pencarian jalur dipicu tepat pada \istilah{state chase}, \istilah{repath} dilakukan secara berkala, dan kinerja tetap terjaga. \\
\addlinespace
Pengujian dan laporan \istilah{playtest} & 20\% &
Terdapat minimal tiga skenario uji, hasil pengamatan yang objektif, dan rumusan perbaikan yang ditindaklanjuti. \\
\addlinespace
Presentasi dan kerja sama & 10\% &
Demonstrasi jelas, setiap anggota memahami bagiannya, serta mampu menjawab pertanyaan penguji. \\
\bottomrule
\end{tabularx}
\end{table}

\subsubsection{Rubrik Proyek Akhir}

\begin{table}[htbp]
\centering
\caption{Rubrik penilaian proyek akhir}
\label{tab:rubrik-proyekakhir}
\begin{tabularx}{\textwidth}{@{}L{4.4cm} C{1.6cm} Y@{}}
\toprule
\headrow \textbf{Kriteria} & \textbf{Bobot} & \textbf{Deskripsi Capaian} \\
\midrule
Fungsionalitas mekanik inti dan level & 30\% &
Level dapat dimainkan dari awal hingga selesai tanpa penghambat; mekanik inti bekerja sesuai rancangan; kontrol pemain responsif. \\
\addlinespace
Integrasi aset visual, animasi, dan UI & 25\% &
Aset visual terpasang rapi, animasi berjalan sesuai keadaan, dan elemen UI menampilkan informasi permainan secara benar. \\
\addlinespace
Integrasi audio dan umpan balik interaksi & 15\% &
BGM dan SFX terpicu pada peristiwa yang tepat dengan volume dan \istilah{timing} yang wajar. \\
\addlinespace
Kualitas kode dan struktur \istilah{project} & 15\% &
\istilah{Project} tersusun rapi, \istilah{script} terorganisasi, penamaan konsisten, dan bebas dari kode mati yang tidak perlu. \\
\addlinespace
\istilah{Playtest}, laporan, dan presentasi & 15\% &
Lembar \istilah{playtest} terisi, \istilah{log bug} lengkap dengan status penanganan, laporan jelas, dan presentasi tepat waktu. \\
\bottomrule
\end{tabularx}
\end{table}

\subsection{Kriteria Ketuntasan dan Konversi Nilai}

\begin{table}[htbp]
\centering
\caption{Konversi nilai akhir}
\label{tab:konversi-nilai}
\begin{tabularx}{\textwidth}{@{}C{2.0cm} C{3.0cm} C{2.4cm} Y@{}}
\toprule
\headrow \textbf{Huruf} & \textbf{Rentang Nilai} & \textbf{Angka Mutu} & \textbf{Kategori Ketercapaian} \\
\midrule
A & 80--100 & 4,00 & Sangat baik; seluruh Sub-CPMK tercapai \\
B & 70--79  & 3,00 & Baik; Sub-CPMK tercapai dengan sedikit kekurangan \\
C & 60--69  & 2,00 & Cukup; Sub-CPMK tercapai pada tingkat minimal \\
D & 50--59  & 1,00 & Kurang; sebagian Sub-CPMK belum tercapai \\
E & $<$ 50  & 0,00 & Gagal; Sub-CPMK tidak tercapai \\
\bottomrule
\end{tabularx}
\tabsource{Konversi nilai mengikuti Pedoman Akademik \NamaPT\ dan dapat disesuaikan apabila terdapat ketentuan berbeda pada program studi.}
\end{table}

Mahasiswa dinyatakan \textbf{lulus} mata kuliah ini apabila memenuhi seluruh
ketentuan berikut.

\begin{enumerate}
  \item Memperoleh nilai akhir minimal \textbf{60} (huruf C).
  \item Memiliki tingkat kehadiran minimal \textbf{75\%} dari total pertemuan.
  \item Menyerahkan seluruh keluaran wajib, yaitu mini proyek AI dan proyek akhir.
\end{enumerate}
